Let $ABC$ be a triangle.
Let $\gamma$ be its incircle.
Let $\alpha_1, \alpha_2, \alpha_3$ be circles such that
$\alpha_1$ is tangent to $\overline{BC}$, to $\overline{CA}$, and to $\gamma$;
$\alpha_2$ is tangent to $\overline{AB}$, to $\overline{CA}$, and to $\gamma$;
$\alpha_3$ is tangent to $\overline{AB}$, to $\overline{BC}$, and to $\gamma$.
Determine the radius of $\gamma$, given the radii of $\alpha_1, \alpha_2, \alpha_3$.

Distinguish all combinations of internal and external tangency between the circles.

\medskip
Alernative formulation.

Three circles are each tangent to a (distinct) unordered pair of (distinct) sides of a triangle
and to the incircle of the triangle.
Determine the radius of the incircle, given the radii of the three circles.

Distinguish all combinations of internal and external tangency between the circles.
